\documentclass[10pt]{beamer}

\beamertemplatenavigationsymbolsempty
\usetheme{Madrid}
\usecolortheme{spruce}
\usepackage[utf8]{inputenc}
\usepackage[english]{babel}
\usepackage{listings} 
\usepackage{xcolor}
\usepackage{graphicx}
\usepackage{pifont}

\lstdefinelanguage{JavaScript}{
  keywords={break, case, catch, continue, debugger, default, delete, do, else, false, finally, for, function, if, in, instanceof, new, null, return, switch, this, throw, true, try, typeof, var, void, while, with, let, const, export, import, await, async},
  morekeywords={play, instrument, stop, setBpm, reverb, pan, ampEnvelope, sequencer, eventList, sub, drums1, synthBass1, drums2},
  sensitive=true,
  morecomment=[l]{//},
  morecomment=[s]{/*}{*/},
  morestring=[b]",
  morestring=[b]'
}

% --- Syntax Highlighting Styles ---
\definecolor{lightgray}{rgb}{0.95, 0.95, 0.95}
\definecolor{darkgray}{rgb}{0.4, 0.4, 0.4}
\definecolor{purple}{rgb}{0.58, 0, 0.82}
\definecolor{blue}{rgb}{0, 0, 1}

\lstset{
    language=JavaScript,
    backgroundcolor=\color{lightgray},
    basicstyle=\ttfamily\scriptsize,
    keywordstyle=\color{blue}\bfseries,
    commentstyle=\color{darkgray}\itshape,
    stringstyle=\color{purple},
    breaklines=true,
    showstringspaces=false,
    captionpos=b
}

% --- Title Information ---
\title[litePlay.js]{litePlay.js: a JavaScript module for \textit{lite coding}}
\subtitle{Concepts, tools, and research results}
\author[Felippe Barros]{Felippe Barros}

\institute[] % (optional)
{
  PhD candidate in Arts - UFES\\
  Federal University of Espírito Santo/Brazil
}
\date{May 6, 2026}

\begin{document}

% Title Slide
\begin{frame}
    \titlepage
\end{frame}

% Table of Contents
\begin{frame}{Table of Contents}
    \tableofcontents
\end{frame}

% --- SECTION 1: INTRODUCTION ---
\section{Introduction to litePlay.js}
\begin{frame}{Lite coding?}
	\begin{itemize}
		\item Ubiquitous music movement
		\item Novel music-making practice featuring musical programming
			with a reduced learning curve
		\item Tailored for non-programmers or casual participants
		\item Exploration of the limits of semantics for musical knowledge-sharing
		\item Seeks to expand participation to musicians with musical
			knowledge that goes beyond the acoustic-instrumental
			and notational logic of Western music
	\end{itemize}
\end{frame}

\begin{frame}{litePlay.js?}
    \begin{itemize}
	\item \textbf{What is it?} A JavaScript module for musical \textit{lite
		coding} designed by Victor Lazzarini
	\item \textbf{Goal:} Prioritize accessibility for beginners and casual
		interaction in musical programming without prescribing aesthetic ideals
        \item \textbf{Technology:} Powered by the Csound WASM (WebAudio) engine
	\item \textbf{Environments:} Online editor\footnote{Available at
		\url{g-ubimus.github.io/litePlay.js}} or integration via a
		    \texttt{<script>} tag
    \end{itemize}
\end{frame}

\begin{frame}[fragile]{Basic syntax}
    The core of the system is the \texttt{play()} function:
    \begin{lstlisting}
// Plays a middle C (C4) with the default piano sound
play(); 

// Set an instrument and play a specific note
instrument(xylophone);
play(E4);

// Stop all sounds
stop();
    \end{lstlisting}
    \vfill
    \textbf{Notes:} Uses midi values or symbolic names from \texttt{Cm1} to
	\texttt{G8} for pitch.
\end{frame}

\begin{frame}{Parameters of an event}
A musical event is defined by five optional attributes:
    \begin{block}{Structure: [what, howLoud, when, howLong, onSomething]}
        \begin{itemize}
            \item \textbf{what:} The note or sound (e.g., \texttt{C4}, \texttt{kick}).
            \item \textbf{howLoud:} Loudness (0 to 1).
            \item \textbf{when:} Start time (seconds).
            \item \textbf{howLong:} Duration (seconds).
            \item \textbf{onSomething:} The instrument (e.g., \texttt{violin}).
        \end{itemize}
    \end{block}
    \begin{exampleblock}{Example}
        \texttt{play([C4, 0.5, 0, 2, banjo])}
    \end{exampleblock}
\end{frame}

\begin{frame}{Generators}
Define a particular scope from which random values will be drawn:
	\begin{block}{Insert a token in place of a value}
        \begin{itemize}
	    \item \textbf{Pitch}: lowPitch, midPitch, highPitch (membranophones, idiophones)
            \item \textbf{Loudness}: soft, midLevel, loud
            \item \textbf{When:} now, soon, later
            \item \textbf{How long:} shortDur, midDur, longDur
	    \item \textbf{On something:} instrument groups (onSynth, onVoice, onBass, etc)
        \end{itemize}
    \end{block}
    \begin{exampleblock}{Example}
        \texttt{play([midPitch, soft, now, shortDur, onPlucked])}
    \end{exampleblock}
\end{frame}

\begin{frame}[fragile]{Sequencing events}
    Works with an automatic grid based on BPM.
    \begin{lstlisting}
setBpm(120);
let backBeat = [kick, snare, kick, snare];
sequencer.add(drums2, backBeat, 0.8, 1); // 1 = Quarter note
sequencer.play();
    \end{lstlisting}
    \textbf{Subdivisions:} \texttt{sub(snare, snare)} divides a beat.

    \begin{lstlisting}
let backBeat = [
	kick, snare, kick, sub(snare, sub(lowTom, tom, hiTom), snare)];
	// O = Silence
    \end{lstlisting}
\end{frame}

\begin{frame}[fragile]{How does it compare to live coding languages?}
    \textbf{Strudel.cc:} uses a highly condensed language, ideal for fast typing.
\begin{lstlisting}[language=JavaScript, backgroundcolor=\color{lightgray}]
$: s("[bd <hh oh>]*2").bank("tr909").dec(.4)
\end{lstlisting}

   \vspace{0.3cm}
    \textbf{litePlay.js:} uses JavaScript structures, prioritizing readability.
\begin{lstlisting}[language=JavaScript, backgroundcolor=\color{lightgray}]
function strudel(){
  let pattern = [bassDrum1, closedHiHat, bassDrum1, openHiHat];
  sequencer.add(drums4, pattern, loud, 0.5);
  sequencer.play();
}

strudel();
\end{lstlisting}
\end{frame}

\section{Study: predicting sounds}
\begin{frame}{Code Comprehension Study}
    \textbf{Hypothesis:} litePlay.js's simplified syntax is more readable for non-programmers than traditional Csound code.
    \begin{block}{Methodology}
        \begin{itemize}
		\item \textbf{Participants:} undergraduate music students with
			minimal experience on programming (n=16, aged from 18
			to 44 years old, avg.= 25.8)
		\item \textbf{Task:} to read code snippets in both programming environments and predict their sonic output.
	    \item \textbf{Metrics:} SEQ (Single Ease Question), Accuracy scores, and SUS (System Usability Scale).
        \end{itemize}
    \end{block}
\end{frame}

\begin{frame}[fragile]{Reading code snippets}
Arguably the simplest Csound instrument:
\begin{lstlisting}
instr Senoide
  iDur = p3
  iAmp = 0.1
  iFreq = 440
  aSine = poscil(iAmp, iFreq)
  outall(aSine)
endin
schedule("Senoide", 0, 5)
\end{lstlisting}
Arpeggios using litePlay.js
\begin{lstlisting}
lp.instrument(lp.vibraphone);
let arpejo = lp.eventList.create(
  [C4, 0.5, 0, 1],
  [E4, 0.7, 1, 1],
  [G4, 0.9, 2, 1]);
let loopArpejo = arpejo.repeat(3, 0);
lp.play(
  [G4, 1, loopArpejo, 2],
  [B4, 1, loopArpejo, 2],
  [D4, 1, loopArpejo, 2]);
    \end{lstlisting}
\end{frame}

\begin{frame}{Results: Readability and Accuracy}
    Data indicated a significant advantage for litePlay.js in semantic readability.
    \begin{table}
        \centering
        \begin{tabular}{|l|c|c|r|}
        \hline
        \textbf{Metric} & \textbf{Csound} & \textbf{litePlay.js} & \textbf{p-value} \\
        \hline
        SEQ (Perceived Ease) & 0.32 & 0.51 & 0.036 \\
        \hline
        Accuracy Score & 0.26 & 0.46 & 0.027 \\
        \hline
	SUS (Usability) & 30 & 38.3 & 0.115 \\
        \hline
        \end{tabular}
    \end{table}
\end{frame}

\begin{frame}{Threats to validity}
\begin{block}{Threats to external validity}
        \begin{itemize}
	    \item \textbf{Population validity}: small, students from the same course, non-English speakers
            \item \textbf{Ecological validity}: generators could give even better readability
        \end{itemize}
    \end{block}
    \begin{block}{Threats to internal validity}
        \begin{itemize}
	    \item \textbf{Construct validity}: other music programming
		    languages could provide different results
	    \item \textbf{Design confound}: IDE design, word choice,
		    indentation, colouring can impact readability and user
			usability
	    \item \textbf{Confounding variable}: a sinewave and an arpeggio are
		    substantially different
        \end{itemize}
    \end{block}
\end{frame}

\begin{frame}{Future perspectives}
\begin{block}{}
        \begin{itemize}
		\item \textbf{Dedicated IDE}: 
			\begin{itemize}
				\item \ding{51} Read-Eval-Print Loop (REPL)
				\item \ding{51} Autocomplete
				\item \ding{51} Audio recording
				\item \ding{55} Audio upload
			\end{itemize}
		\item \textbf{Machine listening}:
			\begin{itemize}
				\item \ding{51} Pitch/percussion, amplitude, duration \ding{55} Timbre, texture...
				\item \ding{55} Agent: sugestions, adaptations (helping with cognitive load)
			\end{itemize}

        \end{itemize}
    \end{block}
\end{frame}

\begin{frame}
    \centering
    \Huge Thank you! \\
    \vspace{1cm}
    \normalsize \texttt{g-ubimus.github.io/litePlay.js/}
    \\
    \normalsize \texttt{g-ubimus.github.io/litePlay.docs/}
\end{frame}

\end{document}
